\subsection{What the comparator signal does and does not supply}
\label{sec:qsp-applicability-audit}

I distinguish a payoff transformation from a mean estimator. The published
payoff construction of \citet{stamatopoulos}, Sections 2--4,
equations (4)--(25), acts on a scalar already held in a binary register.
Its approximation recipe in Appendix A constrains finitely many nodes;
I therefore supply separate admissibility and uniform-error certificates
before transferring that construction to a different financial function.

\begin{proposition}[Exact comparator signal and the identity case]
\label{prop:comparator-qsp-identity}
Let $c\geq0$, $N=2^c$, and let $a$ be a binary integer with
$0\leq a\leq N$. Its register has $c+1$ bits so that the endpoint $N$
is represented. Let $J$ have $c$ qubits, let $F$ be a flag qubit, and
let $C$ be the Boolean permutation
\[
 C\lvert a,j,f\rangle
   =\lvert a,j,f\mathbin\oplus\mathbf1_{\{j<a\}}\rangle.
\]
Define $U=C(I\otimes H^{\otimes c}\otimes I)$ and, within a fixed
$a$ sector, define the input and output projectors
\[
 \Pi=\lvert0^c,0\rangle\langle0^c,0\rvert,
 \qquad \widetilde\Pi=I_J\otimes\lvert1\rangle\langle1\rvert.
\]
Then $\widetilde\Pi U\Pi$ has exactly one non-zero singular value
$s_a=\sqrt{a/N}$ when $a>0$, and is zero when $a=0$.
The odd transformation $P(s)=s$ requires only $U$ and produces
the original comparator probability $a/N$ exactly. Thus an identity
transformation of an already computed finite output supplies no
distinct payoff oracle or estimator.
\end{proposition}

\begin{proof}
The Hadamards prepare $N^{-1/2}\sum_{j=0}^{N-1}\lvert j\rangle$.
The success projector retains precisely the $a$ integers
$j=0,\ldots,a-1$. Therefore
\[
 \widetilde\Pi U\Pi
 =N^{-1/2}\left(\sum_{j=0}^{a-1}\lvert j,1\rangle\right)
       \langle0^c,0\rvert.
\]
The input row vector is normalised. The output vector has squared norm
$a/N$ because its basis vectors are orthonormal. Multiplying its adjoint
by itself gives $(a/N)\Pi$, which proves the singular-value assertion.
The endpoint cases are covered by an empty or complete sum; for $c=0$
the sole integer is $j=0$ and the circuit is a controlled NOT from the
single value bit. The identity transformation is the displayed matrix
itself. It requires no phase approximation, polynomial fitting, or
additional signal calls. An outer preparation over latent values changes
the flag probability to $\sum_y p_y a(y)/N$ by orthogonality of the
retained latent labels, exactly as in the existing finite oracle.
\end{proof}

This identity conclusion depends on the precise signal contract. A
binary value register holding $x$ is not itself a block encoding of $x$;
the comparator above supplies the singular value $\sqrt x$. If a
different upstream oracle instead has singular value $x$, probability
$x$ would require a polynomial approximating $\sqrt x$ on the
nonnegative signal domain, and the identity polynomial would yield
$x^2$. If the upstream signal is a transformed log price, its target
polynomial is a different composition again. The proposition therefore
does not exclude an efficient QSP replacement for a nonlinear operation
that precedes the current numerical output. It identifies the trivial
comparison at the current comparator boundary only.

The comparator convention must be checked rather than copied by name.
The displayed strict predicate $[a<j]$ in equation (18) of the published
paper is inconsistent with the $a$ good zero outcomes displayed in its
equation (19): the literal predicate gives $a+1$ zero outcomes for
$0\leq a<N$. I use $[j<a]$ with good flag one, or equivalently
$[j\geq a]$ with good flag zero. This observation is a direct finite
count, not a claim that the authors' implementation uses the inconsistent
convention. I have not inspected that implementation.

\begin{proposition}[Cost and error specification for a nontrivial transform]
\label{prop:qsp-transform-complete-specification}
Suppose an explicitly retained phase sequence with $d$ signal calls
implements a real polynomial $P$ of definite parity through the
alternating $U,U^\dagger$ construction. Let $h:[0,1]\to[0,1]$ be a
target probability, and suppose
\[
 |P(s)|\leq1\quad(-1\leq s\leq1),\qquad
 |P(\sqrt x)-\sqrt{h(x)}|\leq\eta
 \quad\hbox{for every attainable }x=a/N.
\]
Suppose the emitted phase gates have aggregate operator-norm error at
most $\xi$ relative to that exact phase sequence. Then the implemented
probability differs from $h(x)$ by at most $2\eta+2\xi$ at each attainable
value. The same bound holds for any mixture over the latent labels.
It is valid to replace $2\eta$ by the directly certified maximum of
$|P(\sqrt x)^2-h(x)|$ when that maximum is smaller.
\end{proposition}

\begin{proof}
For $u,v\in[-1,1]$, $|u^2-v^2|=|u-v||u+v|\leq2|u-v|$.
The ideal polynomial probability therefore differs by at most $2\eta$.
For unit vectors $\psi,\varphi$ and any orthogonal projector $E$,
\[
 |\langle\psi,E\psi\rangle-\langle\varphi,E\varphi\rangle|
 \leq |\langle\psi-\varphi,E\psi\rangle|
       +|\langle\varphi,E(\psi-\varphi)\rangle|
 \leq2\|\psi-\varphi\|.
\]
For a product of unitary factors, insertion and subtraction of one
factor at a time gives a telescoping operator-norm bound equal to the
sum of the factor errors. Exact signal calls and exact conjunction
computations add zero. Hence the compiled output state is within $\xi$
of the ideal output, and its projected probability differs by at most
$2\xi$. The triangle inequality gives the stated bound. Taking a
probability-weighted sum preserves it.
\end{proof}

I can implement each signed phase without controlling a measurement
tape. With selector qubit $S$ and a clean bit $B$ computing a projector
predicate $b$, the gate sequence
\[
 \operatorname{CX}_{B,S}\,R_z(2\phi)_S\,
 \operatorname{CX}_{B,S}
 =\exp\{i\phi Z_S(2b-I)\}
\]
follows by conjugating $Z_S$ by the controlled NOT. The formula uses
the convention $R_z(\theta)=\exp(-i\theta Z/2)$; in chronological
order the same three gates are applied. For $\Pi$, $b$ is the
conjunction of the $c+1$ zero tests on $J,F$. A clean chain computes
and uncomputes this predicate using $2c$ exact Toffolis and
$2(c+1)$ NOT gates. For $\widetilde\Pi$, the predicate is the
single flag and requires no conjunction. Each phase therefore includes
two controlled NOT gates and one actual synthesised $R_z$, in addition
to the projector predicate's gates. The selector Hadamards and its
final zero test are retained. In the odd case the good event is
$F=1,S=0$; in the even case it is $J=0^c,F=0,S=0$.
The success reflection must consequently include the full good predicate,
including selector and, for even degree, comparator registers.

If $V$ computes the clean numerical output while retaining the latent
label, the direct complete wrapper is $V^\dagger\mathcal S_PV$ after
latent preparation, where $\mathcal S_P$ is the retained QSP tape.
For any additive gate metric its charge is
\[
 2C(V)+dC(U)+\sum_{k=1}^d C(\text{signed projector phase}_k)
       +2C(H_S),
\]
plus the latent preparation, a compiled success reflection, estimator
calls, observations and reset. $V^\dagger$ is justified because the
transform preserves the value label. A fused compiler can remove
redundant computation, but only an emitted circuit establishes the
corresponding saving. The $P(s)=s$ case can remove the wrapper entirely
and is exactly the existing fused preparation. Non-Clifford phase costs
must be the counts of the retained Clifford+$T$ tapes, not a logarithmic
asymptotic estimate. Their matrix certificates retain relevant global
phases, particularly for later controlled use.

\begin{proposition}[The conditioned residual is not generally a function
of geometric log price alone]
\label{prop:qsp-geometric-scalar-obstruction}
For the continuous financial model let
$m_G(y)=\mu_G+R_G\cdot y$ and
$B_A(y)=\sum_iw_i\exp(\mu_i+R_i\cdot y)$, with $w_i>0$.
Suppose there exists a vector $v$ with $R_G\cdot v=0$ and
$R_i\cdot v\ne0$ for at least one $i$. Then
\[
 h(y)=P_\lambda(\log B_A(y))-P_\lambda(m_G(y))
\]
is not a function of $m_G(y)$ alone, where $P_\lambda$ is the normalised
Gaussian-smoothed capped payoff from this model with $\lambda,H>0$.
\end{proposition}

\begin{proof}
Fix any $y_0$ and parameterise $y=y_0+tv$. The geometric log price is
constant. Differentiating the finite sum twice gives
\[
 \frac{d^2}{dt^2} B_A(y_0+tv)
 =\sum_iw_i e^{\mu_i+R_i\cdot(y_0+tv)}(R_i\cdot v)^2>0.
\]
Every term is nonnegative and the specified term is positive for all
$t$, so $B_A$ is strictly convex and in particular nonconstant on this
line. The logarithm is strictly increasing. The smoothed payoff is
strictly increasing: differentiating its expectation by dominated
convergence gives
\[
 P'_\lambda(m)
 =H^{-1}\mathbb E\left[K e^{m+\lambda Z}
   \mathbf1_{\{K<K e^{m+\lambda Z}<K+H\}}\right]>0.
\]
The integrand is bounded by $(K+H)/H$ and is positive on an interval
with strictly positive Gaussian probability. Thus
$P_\lambda(\log B_A(y))$ is nonconstant on a line on which
$P_\lambda(m_G(y))$ is constant, proving the claim.
\end{proof}

This obstruction applies to the frozen four- and eight-asset loading
matrices, as can be checked using their exact rational decimal entries.
It does not apply to every possible choice of scalar signal, and it does
not assert impossibility of multivariate QSP or a different arithmetic
construction. The two-asset family has one outer coordinate and requires
its own assessment. After $m_A=\log B_A$ has been computed, separate
univariate transforms can evaluate the two smoothed payoffs, but their
combination must preserve the financial difference. In particular,
subtracting square-root amplitudes produces
$(\sqrt{p_A}-\sqrt{p_G})^2$, whereas the required residual is
$p_A-p_G$. These expressions are different whenever
$0<p_G<p_A$. An explicit linear-combination or signed-expectation
construction, its normalisation and its bias budget are therefore needed
before it can replace the positive residual oracle.

The estimator programme of \citet{rallfuller2023} is a separate
comparison. Its final Sections 1.1 and 3 specify rotated projectors and
Chebyshev sampling. I do not transfer empirical query constants into a
worst-case logical-gate budget. I count every confidence look in an
adaptive implementation, including repeated looks at a fixed degree,
and certify the conversion from its amplitude target to the present
probability target.
